\cotefile{1}
\scanpage{108}
% ms. line 01
\uline{Genre d'une fonction entière}

% ms. line 02
Il arrive que l'on ait $\sum\frac{1}{|a_n|^{k+1}}<+\infty$
% ms. line 03
pour quelque $k$, (soit $k$ le plus petit entier qui
% ms. line 04
y satisfasse) et de plus que dans la représentation (2),
% ms. line 05
$g(z)$ soit un polynôme --
% ms. line 06
d'ailleurs déterminé à une constante près.

% ms. line 07
Soit $r$ son degré, le plus \deleted{petit} \inserted{grand} des
% ms. line 08
nombres $r,k$ s'appelle le genre de $f(z)$.
% ms. line 08a
\inserted{C'est un nb \uline{entier} $\geq 0$.}

% ms. line 09
Donc si $f(z)$ est de genre $\leq r$, alors on
% ms. line 10
a la formule
% ms. line 11
\[ f(z)=e^{\alpha_0+\alpha_1 z\ldots+\alpha_{\unsure{n|r}}z^{\unsure{n|r}}}\prod_n\left(1-\frac{z}{a_n}\right)e^{\frac{z}{a_n}+\ldots+\frac{z^r}{r a_n^{\unsure{r}}}}\tag*{(3)} \]

% ms. line 12
Dire que $f(z)$ est de genre zéro, signifie
% ms. line 13
que $f(z)$ est de la forme
% ms. line 14-15
\[ \begin{aligned} \vphantom{f(z)}\\[-2pt]\text{avec}\quad \end{aligned}\left\{\begin{aligned} f(z)&=\mathrm c^{\mathrm{te}}\times\prod_n\left(1-\frac{z}{a_n}\right)\\[-2pt] \sum_n\frac{1}{|a_n|}&<+\infty \end{aligned}\right. \tag*{(4)} \]

% ms. line 16
\uline{Ordre d'une fonction entière} \quad Posons $M(r)=\mathop{\rm Sup}_{z\unsure{=|\leq}r}|f(z)|$
% ms. line 17
Puis $\rho=\overline{\lim}\frac{\log\log M(r)}{\deleted{\unsure{z}}\log r}$ \quad $(0\leq\rho\leq+\infty)$, $\rho$ est
% ms. line 18
dit l'ordre de $f(z)$, et $f(z)$ est dite d'ordre
% ms. line 19
fini si $\rho<+\infty$. Cela signifie qu'il existe
% ms. line 20
$k>0$ tel que $|f(z)|\leq e^{(|z|^k)}$. Alors $\rho$
% ms. line 21
est la \deleted{plus petit nombre} \inserted{borne inférieure des} $\sigma>0$ tels que
% archive sheet number: 107
